\section{Trace norm}

\begin{definition}[Schatten one-norm]\label{def:schatten_one_norm}
    \lean{Matrix.schattenOneNorm}
    \leanok
    For a matrix $A\in\MN{D}$, the singular values $s_0(A),s_1(A),\ldots$ form
    a finitely supported family: only finitely many are nonzero. The Schatten
    one-norm of $A$ is their sum, that is, the sum of the finitely many nonzero
    singular values:
    \[
        \lVert A\rVert_1=\sum_{i\,:\,s_i(A)\neq 0}s_i(A).
    \]
    This is the $p=1$ case of the Schatten $p$-norm
    \cite[Chapter~8, Section~8.1]{Wolf2012Quantum}. The equivalent closed
    formula $\lVert A\rVert_1=\sum_{i=0}^{D-1}s_i(A)$, summing over all $D$
    singular values including trailing zeros, is part of
    Theorem~\ref{thm:trace_norm_elementary}.
\end{definition}

\begin{definition}[Trace norm]\label{def:trace_norm}
    \lean{Matrix.traceNorm}
    \leanok
    \uses{def:schatten_one_norm}
    The trace norm of $A\in\MN{D}$ is its Schatten one-norm:
    \[
        \lVert A\rVert_{\mathrm{tr}}=\lVert A\rVert_1.
    \]
\end{definition}

Wolf records the equivalent formula
$\lVert A\rVert_1=\tr\lvert A\rvert$
\cite[Chapter~8, Section~8.1]{Wolf2012Quantum}; this is
Theorem~\ref{thm:trace_norm_abs}.

\begin{theorem}[Finite singular-value expansions]\label{thm:trace_norm_singular_values}
    \lean{Matrix.schattenOneNorm_eq_sum_support}
    \lean{Matrix.traceNorm_eq_sum_support}
    \lean{Matrix.traceNorm_eq_sum_range_finrank_range}
    \leanok
    \uses{def:schatten_one_norm, def:trace_norm}
    The Schatten one-norm and trace norm are the sums over the finite support of
    the singular-value sequence. Equivalently, the trace norm is the sum over
    the indices below the rank of the represented linear map.
\end{theorem}

\begin{proof}\leanok
    The singular-value sequence satisfies $s_i(A) = 0$ for all
    $i \geq \rank(A)$, so summing over the support and summing
    over $\{0,\ldots,\rank(A)-1\}$ give the same value.
\end{proof}

\begin{theorem}[Elementary trace-norm properties]
    \label{thm:trace_norm_elementary}
    \lean{Matrix.traceNorm_eq_sum_fin}
    \lean{Matrix.schattenOneNorm_nonneg}
    \lean{Matrix.traceNorm_nonneg}
    \lean{Matrix.schattenOneNorm_eq_zero_iff}
    \lean{Matrix.traceNorm_eq_zero_iff}
    \lean{Matrix.traceNorm_zero}
    \lean{Matrix.traceNorm_pos_iff}
    \leanok
    \uses{def:trace_norm, def:schatten_one_norm, thm:trace_norm_singular_values}
    For every $A\in\MN{D}$,
    \[
        \lVert A\rVert_{\mathrm{tr}}
        =\sum_{i=0}^{D-1}s_i(A),\qquad
        \lVert A\rVert_{\mathrm{tr}}\geq 0,
    \]
    and $\lVert A\rVert_{\mathrm{tr}}=0$ if and only if $A=0$.
    Equivalently, $\lVert A\rVert_{\mathrm{tr}}>0$ if and only if $A\neq0$.
\end{theorem}

\begin{proof}\leanok
    The singular values satisfy $s_i(A) \geq 0$, so
    \[
        \lVert A\rVert_{\mathrm{tr}}=0
        \iff \forall i,\; s_i(A)=0
        \iff A=0.
    \]
\end{proof}

\begin{theorem}[Trace norm as trace of the absolute value]
    \label{thm:trace_norm_abs}
    \lean{Matrix.traceNorm_eq_re_trace_abs}
    \leanok
    \uses{def:trace_norm}
    For every $A\in\MN{D}$, with
    $\lvert A\rvert=\sqrt{A^\dagger A}$ the positive-semidefinite square root
    of $A^\dagger A$,
    \[
        \lVert A\rVert_{\mathrm{tr}}
        =\re\bigl(\tr\lvert A\rvert\bigr).
    \]
    This is the formula $\lVert A\rVert_1=\tr[\lvert A\rvert]$
    of \cite[Chapter~8, Section~8.1]{Wolf2012Quantum}; the trace of the
    positive-semidefinite matrix $\lvert A\rvert$ is real.
\end{theorem}

\begin{proof}\leanok
    \uses{lem:trace_norm_sqrt_eigenvalues}
    The map $T^\dagger T$, for $T$ the linear map on $\C^D$ represented by
    $A$, is represented by $A^\dagger A$, so the singular values of $A$ are
    $s_i(A)=\sqrt{\lambda_i}$ with $\lambda_0,\ldots,\lambda_{D-1}$ the
    eigenvalues of $A^\dagger A$. Diagonalizing
    $A^\dagger A=U\diag(\lambda_i)U^\dagger$ gives
    $\lvert A\rvert=U\diag(\sqrt{\lambda_i})U^\dagger$, hence
    \[
        \lVert A\rVert_{\mathrm{tr}}
        =\sum_{i=0}^{D-1}s_i(A)
        =\sum_{i=0}^{D-1}\sqrt{\lambda_i}
        =\tr\lvert A\rvert
        =\re\bigl(\tr\lvert A\rvert\bigr),
    \]
    where the last step uses that $\lvert A\rvert$ is positive semidefinite,
    so $\tr\lvert A\rvert\geq 0$ is real.
\end{proof}

\begin{theorem}[Trace-norm homogeneity]
    \label{thm:trace_norm_homogeneous}
    \lean{Matrix.traceNorm_smul}
    \leanok
    \uses{def:trace_norm}
    For every $c\in\C$ and $A\in\MN{D}$,
    \[
        \lVert cA\rVert_{\mathrm{tr}}
        =\lvert c\rvert\,\lVert A\rVert_{\mathrm{tr}}.
    \]
    This is the homogeneity axiom for matrix norms
    \cite[Chapter~8, Section~8.1]{Wolf2012Quantum}.
\end{theorem}

\begin{proof}\leanok
    \uses{thm:trace_norm_abs}
    From $(cA)^\dagger(cA)=\lvert c\rvert^2\,A^\dagger A$ and uniqueness of
    the positive-semidefinite square root,
    $\lvert cA\rvert=\lvert c\rvert\,\lvert A\rvert$. Linearity of the trace
    gives
    $\tr\lvert cA\rvert
    =\lvert c\rvert\,\tr\lvert A\rvert$, and the claim follows
    from Theorem~\ref{thm:trace_norm_abs}.
\end{proof}

\begin{theorem}[Trace-norm unitary invariance]
    \label{thm:trace_norm_unitary_invariance}
    \lean{Matrix.traceNorm_unitary_mul}
    \lean{Matrix.traceNorm_mul_unitary}
    \lean{Matrix.traceNorm_unitary_mul_unitary}
    \leanok
    \uses{def:trace_norm}
    For all unitaries $U,V\in\MN{D}$ and every $A\in\MN{D}$,
    \[
        \lVert UA\rVert_{\mathrm{tr}}
        =\lVert AV\rVert_{\mathrm{tr}}
        =\lVert UAV\rVert_{\mathrm{tr}}
        =\lVert A\rVert_{\mathrm{tr}}.
    \]
    The trace norm is thus unitarily invariant
    \cite[Chapter~8, Section~8.1]{Wolf2012Quantum}.
\end{theorem}

\begin{proof}\leanok
    \uses{thm:trace_norm_abs}
    From $(UA)^\dagger(UA)=A^\dagger U^\dagger UA=A^\dagger A$ the absolute
    values agree: $\lvert UA\rvert=\lvert A\rvert$. For the right factor,
    $(AV)^\dagger(AV)=V^\dagger(A^\dagger A)V$, and since
    $V^\dagger\lvert A\rvert V$ is positive semidefinite with
    \[
        \bigl(V^\dagger\lvert A\rvert V\bigr)^2
        =V^\dagger\lvert A\rvert\,VV^\dagger\,\lvert A\rvert V
        =V^\dagger(A^\dagger A)V,
    \]
    uniqueness of the positive-semidefinite square root gives
    $\lvert AV\rvert=V^\dagger\lvert A\rvert V$. Cyclicity of the trace then
    yields
    $\tr\lvert AV\rvert
    =\tr(\lvert A\rvert\,VV^\dagger)
    =\tr\lvert A\rvert$. For the two-sided case,
    $\lVert UAV\rVert_{\mathrm{tr}}=\lVert UA\rVert_{\mathrm{tr}}=\lVert A\rVert_{\mathrm{tr}}$
    by applying the right and left cases in turn.
\end{proof}

\begin{lemma}[Trace norm from the spectrum of $A^\dagger A$]
    \label{lem:trace_norm_sqrt_eigenvalues}
    \lean{Matrix.traceNorm_eq_sum_sqrt_eigenvalues_adjoint_comp_self}
    \leanok
    \uses{def:trace_norm}
    For every $A\in\MN{D}$, the trace norm is the sum of the square roots of
    the eigenvalues of the positive operator $A^\dagger A$:
    \[
        \lVert A\rVert_{\mathrm{tr}}
        = \sum_{i=0}^{D-1}\sqrt{\lambda_i(A^\dagger A)}.
    \]
\end{lemma}

\begin{proof}\leanok
    \uses{thm:trace_norm_elementary}
    Since $s_i(A)=\sqrt{\lambda_i(A^\dagger A)}$ by the definition of singular
    value, the finite singular-value expansion of
    Theorem~\ref{thm:trace_norm_elementary} gives
    \[
        \lVert A\rVert_{\mathrm{tr}}
        = \sum_{i=0}^{D-1}s_i(A)
        = \sum_{i=0}^{D-1}\sqrt{\lambda_i(A^\dagger A)}.
    \]
\end{proof}

\begin{lemma}[Trace norm from the matrix eigenvalues of $A^\dagger A$]
    \label{lem:trace_norm_sqrt_matrix_eigenvalues}
    \lean{Matrix.traceNorm_eq_sum_sqrt_eigenvalues}
    \leanok
    \uses{def:trace_norm}
    For every $A\in\MN{D}$, with $\lambda_0,\ldots,\lambda_{D-1}$ the
    eigenvalues of the matrix $A^\dagger A$,
    \[
        \lVert A\rVert_{\mathrm{tr}}=\sum_{i=0}^{D-1}\sqrt{\lambda_i}.
    \]
\end{lemma}

\begin{proof}\leanok
    \uses{thm:trace_norm_abs}
    Diagonalizing $A^\dagger A=V\diag(\lambda_i)V^\dagger$ gives
    $\lvert A\rvert=V\diag(\sqrt{\lambda_i})V^\dagger$, so
    \[
        \tr\lvert A\rvert=\sum_{i=0}^{D-1}\sqrt{\lambda_i},
    \]
    and the claim follows from Theorem~\ref{thm:trace_norm_abs}.
\end{proof}

\begin{theorem}[Variational unitary formula for the trace norm]
    \label{thm:trace_norm_variational}
    \lean{Matrix.norm_trace_conjTranspose_mul_unitary_le}
    \lean{Matrix.exists_mem_unitaryGroup_trace_conjTranspose_mul_eq}
    \lean{Matrix.isGreatest_norm_trace_conjTranspose_mul_unitary}
    \lean{Matrix.traceNorm_eq_sSup_norm_trace_conjTranspose_mul_unitary}
    \leanok
    \uses{def:trace_norm}
    For every $A\in\MN{D}$,
    \[
        \lVert A\rVert_{\mathrm{tr}}
        =\max\left\{\bigl|\tr[A^\dagger U]\bigr|
        \;\middle|\; UU^\dagger=\Id\right\},
    \]
    and the maximum is attained: some unitary $U$ satisfies
    $\tr[A^\dagger U]=\lVert A\rVert_{\mathrm{tr}}$.
    See \cite[Chapter~8, Eq.~(8.11)]{Wolf2012Quantum}.
\end{theorem}

\begin{proof}\leanok
    \uses{lem:trace_norm_sqrt_matrix_eigenvalues}
    For the upper bound, let $v_0,\ldots,v_{D-1}$ be an orthonormal
    eigenbasis of $A^\dagger A$ with eigenvalues $\lambda_i$, and let $U$ be
    unitary. Expanding the trace in this basis and applying the
    Cauchy--Schwarz inequality,
    \[
        \bigl|\tr[A^\dagger U]\bigr|
        =\Bigl|\sum_{i}\langle Av_i,\,Uv_i\rangle\Bigr|
        \le\sum_{i}\lVert Av_i\rVert\,\lVert Uv_i\rVert
        =\sum_{i}\sqrt{\lambda_i}
        =\lVert A\rVert_{\mathrm{tr}},
    \]
    since $\lVert Av_i\rVert^2=\langle v_i,\,A^\dagger A\,v_i\rangle
    =\lambda_i$ and $\lVert Uv_i\rVert=1$.

    For attainment, the vectors $w_i=Av_i/\sqrt{\lambda_i}$, taken over the
    indices with $\lambda_i\neq0$, satisfy
    \[
        \langle w_i,w_j\rangle
        =\frac{\langle v_i,\,A^\dagger A\,v_j\rangle}
              {\sqrt{\lambda_i}\sqrt{\lambda_j}}
        =\delta_{ij},
    \]
    so they form an orthonormal family. Extend it to an orthonormal basis
    $(w_i)_{i}$ of $\C^D$ and let $U$ be the unitary with $Uv_i=w_i$ for all
    $i$. Then
    \[
        \tr[A^\dagger U]
        =\sum_{i}\langle Av_i,\,w_i\rangle
        =\sum_{\lambda_i\neq0}\frac{\lVert Av_i\rVert^2}{\sqrt{\lambda_i}}
        =\sum_{i}\sqrt{\lambda_i}
        =\lVert A\rVert_{\mathrm{tr}},
    \]
    where the indices with $\lambda_i=0$ contribute $0$ to both sides because
    $\lVert Av_i\rVert^2=\lambda_i=0$ forces $Av_i=0$.
\end{proof}

\begin{theorem}[Trace-norm triangle inequality]
    \label{thm:trace_norm_triangle}
    \lean{Matrix.traceNorm_add_le}
    \leanok
    \uses{def:trace_norm}
    For all $A,B\in\MN{D}$,
    \[
        \lVert A+B\rVert_{\mathrm{tr}}
        \le\lVert A\rVert_{\mathrm{tr}}+\lVert B\rVert_{\mathrm{tr}}.
    \]
    Together with homogeneity (Theorem~\ref{thm:trace_norm_homogeneous}) and
    definiteness (Theorem~\ref{thm:trace_norm_elementary}), this completes
    the norm axioms of \cite[Chapter~8, Section~8.1]{Wolf2012Quantum} for the
    trace norm.
\end{theorem}

\begin{proof}\leanok
    \uses{thm:trace_norm_variational}
    Choose by Theorem~\ref{thm:trace_norm_variational} a unitary $U$ with
    $\tr[(A+B)^\dagger U]=\lVert A+B\rVert_{\mathrm{tr}}$.
    Splitting the trace and applying the upper-bound half of the same
    theorem to $A$ and to $B$ separately,
    \[
        \lVert A+B\rVert_{\mathrm{tr}}
        =\bigl|\tr[A^\dagger U]
              +\tr[B^\dagger U]\bigr|
        \le\bigl|\tr[A^\dagger U]\bigr|
          +\bigl|\tr[B^\dagger U]\bigr|
        \le\lVert A\rVert_{\mathrm{tr}}+\lVert B\rVert_{\mathrm{tr}}.
    \]
\end{proof}

\begin{lemma}[Jordan trace-norm formula]
    \label{lem:trace_norm_jordan}
    \lean{Matrix.IsHermitian.traceNorm_eq_re_trace_posPart_add_negPart}
    \leanok
    \uses{def:trace_norm}
    Let $H\in\MN{D}$ be Hermitian with Jordan decomposition $H=H^+-H^-$
    into orthogonal positive parts, $H^\pm\ge0$ and $H^+H^-=0$. Then
    \[
        \lVert H\rVert_{\mathrm{tr}}
        =\tr[H^+]+\tr[H^-].
    \]
\end{lemma}

\begin{proof}\leanok
    \uses{thm:trace_norm_abs}
    The sum $H^++H^-$ is positive semidefinite, and by orthogonality its
    square is
    \[
        (H^++H^-)^2=(H^+)^2+(H^-)^2=(H^+-H^-)^2=H^\dagger H,
    \]
    so $H^++H^-=\sqrt{H^\dagger H}=|H|$. Hence
    \[
        \lVert H\rVert_{\mathrm{tr}}
        =\tr|H|
        =\tr[H^+]+\tr[H^-]
    \]
    by Theorem~\ref{thm:trace_norm_abs}.
\end{proof}

\begin{lemma}[Support projection of the positive part]
    \label{lem:positive_part_support_projection}
    \lean{Matrix.IsHermitian.posPartSupportProj}
    \lean{Matrix.IsHermitian.posPartSupportProj_posSemidef}
    \lean{Matrix.IsHermitian.posPartSupportProj_le_one}
    \lean{Matrix.IsHermitian.posPartSupportProj_mul_posPartSupportProj}
    \lean{Matrix.IsHermitian.posPartSupportProj_mul_self}
    \leanok
    Let $A\in\MN{D}$ be Hermitian with positive part $A^+$. There is a
    matrix $\Pi$ with $0\le\Pi\le\Id$, $\Pi^2=\Pi$, and $\Pi A=A^+$, namely
    the orthogonal projection onto the support space of $A^+$.
\end{lemma}

\begin{proof}\leanok
    Diagonalize
    $A=U\diag(\lambda_1,\dots,\lambda_D)\,U^\dagger$ and set
    \[
        \Pi=U\diag
        \bigl(\chi(\lambda_1),\dots,\chi(\lambda_D)\bigr)\,U^\dagger,
    \]
    where $\chi$ is the indicator function of $(0,\infty)$. Each claim is
    read off eigenvalue-wise: $0\le\chi\le1$ gives $0\le\Pi\le\Id$,
    $\chi^2=\chi$ gives $\Pi^2=\Pi$, and
    $\chi(\lambda)\,\lambda=\max(\lambda,0)$ gives $\Pi A=A^+$.
\end{proof}

\begin{lemma}[Positive-semidefinite trace pairing bounds]
    \label{lem:psd_trace_pairing_bounds}
    \lean{Matrix.PosSemidef.re_trace_mul_nonneg}
    \lean{Matrix.PosSemidef.re_trace_mul_le_of_le_one}
    \leanok
    Let $X,C\in\MN{D}$ with $X\ge0$. If $C\ge0$, then
    $\tr[CX]\ge0$; if $C\le\Id$, then
    $\tr[CX]\le\tr[X]$.
\end{lemma}

\begin{proof}\leanok
    \uses{lem:psd_trace_product_nonneg}
    The first bound is Lemma~\ref{lem:psd_trace_product_nonneg}. For the
    second,
    \[
        \tr[X]-\tr[CX]
        =\tr[(\Id-C)X]\ge0
    \]
    by the same lemma, since $\Id-C\ge0$.
\end{proof}

\begin{lemma}[Positive maps preserve hermiticity]
    \label{lem:positive_map_hermitian}
    \lean{Matrix.isHermitian_map_of_positive}
    \leanok
    \uses{def:positive_map_rectangular}
    If $T:\MN{D}\to\MN{D'}$ is a positive linear map and $H\in\MN{D}$ is
    Hermitian, then $T(H)$ is Hermitian.
\end{lemma}

\begin{proof}\leanok
    \uses{def:positive_map_rectangular}
    Write $H=H^+-H^-$ with $H^\pm\ge0$. Then $T(H)=T(H^+)-T(H^-)$ is a
    difference of positive semidefinite matrices, hence Hermitian.
\end{proof}

\begin{lemma}[Positive-part trace estimate]
    \label{lem:positive_part_trace_estimate}
    \lean{Matrix.re_trace_posPart_map_le}
    \leanok
    \uses{def:positive_map_rectangular, def:trace_preserving_rectangular}
    Let $T:\MN{D}\to\MN{D'}$ be a trace-preserving positive linear map and
    let $H\in\MN{D}$ be Hermitian, with Jordan decompositions $H=P_+-P_-$
    and $T(H)=Q_+-Q_-$. Then
    \[
        \tr[Q_+]\le\tr[P_+].
    \]
\end{lemma}

\begin{proof}\leanok
    \uses{lem:positive_part_support_projection,
      lem:psd_trace_pairing_bounds, lem:positive_map_hermitian,
      def:trace_preserving_rectangular}
    The matrix $T(H)$ is Hermitian by
    Lemma~\ref{lem:positive_map_hermitian}. Let $\Pi_+$ be the support
    projection of $Q_+$ from
    Lemma~\ref{lem:positive_part_support_projection}, so that
    $\Pi_+T(H)=Q_+$ and $0\le\Pi_+\le\Id$. Then
    \[
        \tr[Q_+]
        =\tr[\Pi_+T(P_+)]-\tr[\Pi_+T(P_-)]
        \le\tr[\Pi_+T(P_+)]
        \le\tr[T(P_+)]
        =\tr[P_+],
    \]
    where the first inequality uses $\tr[\Pi_+T(P_-)]\ge0$
    and the second uses $\Pi_+\le\Id$ together with $T(P_+)\ge0$
    (Lemma~\ref{lem:psd_trace_pairing_bounds}); the final step is trace
    preservation.
\end{proof}

\begin{theorem}[Trace-norm contractivity]
    \label{thm:trace_norm_contractivity}
    \lean{Matrix.traceNorm_map_le_of_positive_of_tracePreserving}
    \leanok
    \uses{def:trace_norm, def:positive_map_rectangular,
      def:trace_preserving_rectangular}
    Let $T:\MN{D}\to\MN{D'}$ be a trace-preserving positive linear map.
    Then for all Hermitian $H\in\MN{D}$,
    \[
        \lVert T(H)\rVert_{\mathrm{tr}}\le\lVert H\rVert_{\mathrm{tr}}.
    \]
    See \cite[Chapter~8, Theorem~8.16]{Wolf2012Quantum}.
\end{theorem}

\begin{proof}\leanok
    \uses{lem:trace_norm_jordan, lem:positive_part_trace_estimate,
      lem:positive_map_hermitian}
    Write $H=P_+-P_-$ and $T(H)=Q_+-Q_-$ for the Jordan decompositions.
    Applying Lemma~\ref{lem:positive_part_trace_estimate} to $H$ and to
    $-H$ gives $\tr[Q_+]\le\tr[P_+]$ and
    $\tr[Q_-]\le\tr[P_-]$, so by the Jordan
    trace-norm formula (Lemma~\ref{lem:trace_norm_jordan}),
    \[
        \lVert T(H)\rVert_{\mathrm{tr}}
        =\tr[Q_+]+\tr[Q_-]
        \le\tr[P_+]+\tr[P_-]
        =\lVert H\rVert_{\mathrm{tr}}.
    \]
\end{proof}

\begin{corollary}[Trace-norm contractivity on states]
    \label{cor:trace_norm_contractivity_states}
    \lean{Matrix.traceNorm_map_sub_map_le_of_positive_of_tracePreserving}
    \leanok
    \uses{def:trace_norm, def:positive_map_rectangular,
      def:trace_preserving_rectangular, def:density_matrices}
    Let $T:\MN{D}\to\MN{D'}$ be a trace-preserving positive linear map.
    Then for all density matrices $\rho_1,\rho_2\in\MN{D}$,
    \[
        \lVert T(\rho_1)-T(\rho_2)\rVert_{\mathrm{tr}}
        \le\lVert\rho_1-\rho_2\rVert_{\mathrm{tr}}.
    \]
    See \cite[Chapter~8, Eq.~(8.80)]{Wolf2012Quantum}.
\end{corollary}

\begin{proof}\leanok
    \uses{thm:trace_norm_contractivity}
    By linearity $T(\rho_1)-T(\rho_2)=T(\rho_1-\rho_2)$, and
    $\rho_1-\rho_2$ is Hermitian as a difference of positive semidefinite
    matrices, so Theorem~\ref{thm:trace_norm_contractivity} applies.
\end{proof}
